% =====================================================================
%  Section 5 -- Results
% =====================================================================

\section{Results}
\label{sec:results}

This section is organised around Tables~\ref{tab:t1-footprints}--\ref{tab:t5-latency} and Figures~\ref{fig:f2-pareto}--\ref{fig:f4-throughput}. The pipeline schematic (Figure~\ref{fig:f1-pipeline}) and qualitative panel (Figure~\ref{fig:f3-qualitative}) appear in the supplementary material.

\subsection{Teacher Performance and Karc{\i} \emph{et al.}~(2026) Comparison}
\label{sec:results-teacher}

Our UNet $+$ EfficientNet-B0 teacher achieves \textbf{IoU $= \teacherIoU \pm \teacherStd$} on FSSD across the three folds (F1 $= 0.9611 \pm 0.0044$, accuracy $= 0.9729 \pm 0.0025$). The reference UNet $+$ EfficientNet-B7 number reported by Karc{\i} \emph{et al.}~\cite{karci2026deep} is IoU $= 0.927$, achieved with approximately $10\times$ the parameters ($\sim 63$\,M vs.~our 6.25\,M). The 0.0013 IoU gap between our B0 teacher and the published B7 reference is within the cross-fold standard deviation of both studies and indicates that the FSSD benchmark saturates well before the B7 capacity is required; B0 is sufficient as the teacher.

\subsection{Students Match or Exceed the Teacher Without Distillation (Negative-KD Finding)}
\label{sec:results-kd}

The full segmentation results across all 14 configurations $\times$ 3 folds appear in Table~\ref{tab:t2-segmentation}. Three observations:

% =====================================================================
%  Table T2 -- In-Distribution Segmentation
%  Source of truth: code/runs/3fold/summary.json
% =====================================================================
\begin{table*}[ht]
\centering
\caption{\textbf{In-distribution segmentation performance.} Mean $\pm$ std over 3 folds (FSSD val, $256{\times}256$). Best per metric in \textbf{bold}.}
\label{tab:t2-segmentation}
\small
\setlength{\tabcolsep}{4pt}
\begin{tabular}{lccccc}
\toprule
\textbf{Model / KD}                & \textbf{Acc}             & \textbf{Prec}            & \textbf{Rec}             & \textbf{F1}              & \textbf{IoU}             \\
\midrule
Teacher (UNet$+$EffB0)             & 0.9729 $\pm$ 0.0025      & 0.9623 $\pm$ 0.0077      & 0.9604 $\pm$ 0.0010      & 0.9611 $\pm$ 0.0044      & 0.9257 $\pm$ 0.0078      \\
MobileNetV2 baseline               & 0.9729 $\pm$ 0.0010      & 0.9632 $\pm$ 0.0052      & 0.9595 $\pm$ 0.0028      & 0.9611 $\pm$ 0.0019      & 0.9257 $\pm$ 0.0032      \\
\midrule
MobileNetV3-S / none               & 0.9732 $\pm$ 0.0006      & 0.9620 $\pm$ 0.0068      & 0.9620 $\pm$ 0.0050      & 0.9618 $\pm$ 0.0024      & 0.9269 $\pm$ 0.0043      \\
MobileNetV3-S / resp               & 0.9719 $\pm$ 0.0007      & 0.9599 $\pm$ 0.0050      & 0.9598 $\pm$ 0.0049      & 0.9596 $\pm$ 0.0016      & 0.9229 $\pm$ 0.0028      \\
MobileNetV3-S / feat               & 0.9732 $\pm$ 0.0012      & 0.9616 $\pm$ 0.0040      & 0.9626 $\pm$ 0.0025      & 0.9619 $\pm$ 0.0020      & 0.9271 $\pm$ 0.0035      \\
MobileNetV3-S / comb               & 0.9723 $\pm$ 0.0019      & 0.9603 $\pm$ 0.0062      & 0.9608 $\pm$ 0.0024      & 0.9604 $\pm$ 0.0034      & 0.9242 $\pm$ 0.0061      \\
\midrule
EfficientNet-Lite0 / none          & 0.9735 $\pm$ 0.0012      & 0.9654 $\pm$ 0.0048      & 0.9595 $\pm$ 0.0034      & 0.9622 $\pm$ 0.0026      & 0.9277 $\pm$ 0.0047      \\
EfficientNet-Lite0 / resp          & 0.9729 $\pm$ 0.0013      & 0.9622 $\pm$ 0.0037      & 0.9601 $\pm$ 0.0028      & 0.9609 $\pm$ 0.0007      & 0.9253 $\pm$ 0.0011      \\
EfficientNet-Lite0 / feat          & 0.9734 $\pm$ 0.0008      & 0.9654 $\pm$ 0.0054      & 0.9593 $\pm$ 0.0044      & 0.9621 $\pm$ 0.0020      & 0.9275 $\pm$ 0.0034      \\
EfficientNet-Lite0 / comb          & 0.9733 $\pm$ 0.0016      & 0.9633 $\pm$ 0.0051      & 0.9601 $\pm$ 0.0019      & 0.9615 $\pm$ 0.0021      & 0.9264 $\pm$ 0.0037      \\
\midrule
MobileViT-XXS / none               & \textbf{0.9771 $\pm$ 0.0023} & \textbf{0.9695 $\pm$ 0.0040} & \textbf{0.9650 $\pm$ 0.0025} & \textbf{0.9671 $\pm$ 0.0031} & \textbf{0.9367 $\pm$ 0.0055} \\
MobileViT-XXS / resp               & 0.9757 $\pm$ 0.0019      & 0.9659 $\pm$ 0.0055      & 0.9642 $\pm$ 0.0016      & 0.9649 $\pm$ 0.0033      & 0.9326 $\pm$ 0.0060      \\
MobileViT-XXS / feat               & 0.9766 $\pm$ 0.0014      & 0.9680 $\pm$ 0.0051      & 0.9655 $\pm$ 0.0039      & 0.9665 $\pm$ 0.0026      & 0.9357 $\pm$ 0.0046      \\
MobileViT-XXS / comb               & 0.9752 $\pm$ 0.0025      & 0.9657 $\pm$ 0.0053      & 0.9636 $\pm$ 0.0025      & 0.9644 $\pm$ 0.0039      & 0.9318 $\pm$ 0.0070      \\
\bottomrule
\end{tabular}
\end{table*}


\begin{enumerate}
\item \textbf{The best student exceeds the teacher.} MobileViT-XXS in the no-KD configuration attains IoU $= \bestStudentIoU \pm \bestStudentStd$, \emph{0.011 IoU points above the teacher} (\teacherIoU $\pm$ \teacherStd). The student is also $2.0\times$ smaller in parameters (3.09\,M vs.~6.25\,M). This is the headline result against which the original KD-recipe framing of this paper has to be evaluated.
\item \textbf{EfficientNet-Lite0 (no-KD) matches the teacher within standard deviation} ($0.9277 \pm 0.0047$ vs.~$\teacherIoU \pm \teacherStd$). MobileNetV3-Small (no-KD) likewise matches it ($0.9269 \pm 0.0043$).
\item \textbf{The MobileNetV2 baseline also matches the teacher} ($0.9257 \pm 0.0032$). The teacher's parameter count is therefore not the source of any advantage; it is an over-parameterised lightweight UNet, not a ``heavy'' model in the regime where KD typically helps.
\end{enumerate}

We interpret these observations as direct evidence that FSSD, at 663 RGB images and IoU $\approx 0.93$, is saturated for the ImageNet-pretrained UNet family. The variance across $(\text{model} \times \text{KD})$ cells of Table~\ref{tab:t2-segmentation} is dominated by fold-level seed noise rather than by genuine architectural separations. KD has no headroom to transfer.

Table~\ref{tab:t3-kd-ablation} isolates the KD ablation on the best student. None of the three KD configurations improves on the no-KD baseline:

% =====================================================================
%  Table T3 -- KD Ablation on Best Student
%  Source of truth: code/runs/3fold/summary.json + stats.json (Wilcoxon)
% =====================================================================
\begin{table}[ht]
\centering
\caption{\textbf{KD ablation on the best student (MobileViT-XXS).} Wilcoxon two-sided; $n=3$ folds gives a minimum reportable $p \approx 0.25$.}
\label{tab:t3-kd-ablation}
\small
\begin{tabular}{lccc}
\toprule
\textbf{KD config} & \textbf{IoU (mean $\pm$ std)} & \textbf{$\Delta$ IoU vs.\ none} & \textbf{Wilcoxon $p$} \\
\midrule
none      & 0.9367 $\pm$ 0.0055 & ---      & --- \\
resp      & 0.9326 $\pm$ 0.0060 & $-0.0041$ & n/a \\
feat      & 0.9357 $\pm$ 0.0046 & $-0.0011$ & n/a \\
comb      & 0.9318 $\pm$ 0.0070 & $-0.0049$ & 0.250 \\
\bottomrule
\end{tabular}
\end{table}


The Wilcoxon test on no-KD vs.~comb-KD yields $p = 0.250$ --- the minimum two-sided p-value reportable at $n=3$ --- which means the test is consistent with both ``comb-KD is no worse'' and ``comb-KD is meaningfully worse'' under low power, but in no case does it support ``comb-KD is better''. The remaining contrasts (resp, feat) cannot be tested due to insufficient direction agreement across only three folds. The same qualitative pattern holds on MobileNetV3-Small and EfficientNet-Lite0 (see full Table~\ref{tab:t2-segmentation}). Across \emph{all three architectures $\times$ three KD configurations $\times$ three folds} (27 fold-level data points), the combined KD configuration produces a higher mean IoU than the no-KD baseline in zero out of three architectures.

\subsection{INT8 Quantization is Architecture-Dependent}
\label{sec:results-quant}

Table~\ref{tab:t4-quant} summarises post-training INT8 quantization for the three deployable students:

% =====================================================================
%  Table T4 -- Quantization
%  Source of truth: code/runs/3fold/quantized.json + ONNX file sizes.
% =====================================================================
\begin{table*}[ht]
\centering
\caption{\textbf{Post-training INT8 quantization} (static QDQ, Percentile-99.999\% calibration, 200 train images). EfficientNet-Lite0 (highlighted) is the recommended deployable configuration; the other two students suffer architecture-driven PTQ collapse (Section~\ref{sec:results-quant}).}
\label{tab:t4-quant}
\small
\setlength{\tabcolsep}{5pt}
\begin{tabular}{lccccccc}
\toprule
\textbf{Model} & \textbf{FP32 IoU} & \textbf{INT8 IoU} & \textbf{$\Delta$ IoU} & \textbf{FP32 ONNX (MB)} & \textbf{INT8 ONNX (MB)} & \textbf{Size reduction} \\
\midrule
MobileNetV3-Small   & 0.9269 $\pm$ 0.0043 & 0.4762 $\pm$ 0.3399 & $-0.4507$ & 13.8 & 3.7 & $3.69\times$ \\
\textbf{EfficientNet-Lite0} & \textbf{0.9277 $\pm$ 0.0047} & \textbf{0.8851 $\pm$ 0.0149} & $\boldsymbol{-0.0426}$ & \textbf{19.8} & \textbf{5.2} & $\boldsymbol{3.83\times}$ \\
MobileViT-XXS       & 0.9367 $\pm$ 0.0055 & 0.4752 $\pm$ 0.1936 & $-0.4616$ & 12.0 & 3.5 & $3.42\times$ \\
\bottomrule
\end{tabular}
\end{table*}


The cross-architecture spread is the central finding of this section. Under an identical Percentile-99.999\% static QDQ recipe, EfficientNet-Lite0 retains 95\% of its FP32 IoU with a tight $\pm 0.015$ cross-fold standard deviation, whereas MobileNetV3-Small and MobileViT-XXS both lose $\sim 0.45$ IoU points with cross-fold standard deviations of 0.19--0.34 --- an order of magnitude larger than EfficientNet-Lite0's. We interpret these failures architecturally:

\begin{itemize}
\item \textbf{MobileNetV3-Small} uses the HardSwish activation $x \cdot \mathrm{ReLU6}(x+3)/6$, which has an unbounded positive tail and produces calibration statistics dominated by single-image outliers. Per-tensor INT8 calibration cannot fit a single scale that simultaneously preserves the activations of common flood-pixel patches and of bright reflective outliers (wet asphalt, water-surface sun-glare). Percentile-99.999\% mitigates but does not eliminate this; we conjecture that MobileNetV3-S requires either (a)~QAT to learn quantizer-friendly activations or (b)~a HardSwish\,$\rightarrow$\,ReLU6 surgical replacement before quantization.
\item \textbf{MobileViT-XXS} combines a CNN trunk with self-attention blocks containing LayerNorm. The LayerNorm $\gamma / \beta$ tensors are rank-1 across channels, which means per-channel weight quantization (the standard recipe for Conv layers) is ill-defined; restricting per-channel quantization to \texttt{Conv/Gemm/MatMul} (Section~\ref{sec:methods-ptq}) partially fixes this but leaves the attention \texttt{MatMul} itself quantized per-tensor, where the high dynamic range of attention scores again exceeds what static calibration can fit.
\item \textbf{EfficientNet-Lite0} uses ReLU6 throughout. ReLU6 is bounded by construction at 6.0, so the per-tensor activation calibration finds a stable scale that matches the actual activation range in every fold. The 0.043 IoU drop is consistent with the conventional-wisdom ``$\sim 1$ IoU point'' PTQ tax in the classification literature.
\end{itemize}

The practical implication: \emph{bounded-activation backbones (ReLU6) are the safe default for INT8 PTQ on flood segmentation; HardSwish and attention-based backbones currently require QAT or remain FP32 in deployment.} We do not attempt QAT in this paper --- it is out of scope and an order-of-magnitude larger training investment --- and instead recommend EfficientNet-Lite0 as the deployable configuration on the basis of its joint accuracy $+$ size $+$ stability behaviour.

\subsection{Per-Platform Inference Latency}
\label{sec:results-latency}

Table~\ref{tab:t5-latency} reports inference latency on the two deployment platforms.

% =====================================================================
%  Table T5 -- Per-Platform Latency
%  Source of truth: code/runs/3fold/{lat_m2,lat_wasm}.json
% =====================================================================
\begin{table*}[ht]
\centering
\caption{\textbf{Per-platform inference latency.} Single-image batches, $256{\times}256$ input. M2~Pro and browser WASM both consume the same \texttt{.onnx} artifact under different execution providers. INT8 WASM rows are n/a because ONNX Runtime Web 1.17 does not implement QDQ-format INT8 Conv (upstream limitation; scheduled for ORT-Web 1.18$+$).}
\label{tab:t5-latency}
\small
\setlength{\tabcolsep}{5pt}
\begin{tabular}{lcccccc}
\toprule
\multirow{2}{*}{\textbf{Model}}
  & \multicolumn{3}{c}{\textbf{Apple M2~Pro (ORT-CPU)}}
  & \multicolumn{3}{c}{\textbf{Browser WASM SIMD (ORT-Web)}} \\
\cmidrule(lr){2-4} \cmidrule(lr){5-7}
  & p50 (ms) & p95 (ms) & FPS  & p50 (ms) & p95 (ms) & FPS  \\
\midrule
Teacher (UNet$+$EffB0) FP32        & 274.0 & 280.7 & 3.6   & \multicolumn{3}{c}{n/a (out-of-scope for browser deploy)} \\
\midrule
MobileNetV3-Small FP32             &  19.5 &  20.4 & 51.2  & 133.3 & 134.8 & 7.5 \\
MobileNetV3-Small INT8             &  13.7 &  14.1 & 72.9  & \multicolumn{3}{c}{n/a (ORT-Web 1.17 limitation)} \\
EfficientNet-Lite0 FP32            &  34.2 &  52.7 & 29.3  & 163.8 & 210.5 & 6.1 \\
\textbf{EfficientNet-Lite0 INT8}   &  \textbf{17.3} & \textbf{20.1} & \textbf{57.7} & \multicolumn{3}{c}{n/a (ORT-Web 1.17 limitation)} \\
MobileViT-XXS FP32                 &  28.3 &  32.8 & 35.4  & 149.4 & 159.8 & 6.7 \\
MobileViT-XXS INT8                 &  21.6 &  25.0 & 46.3  & \multicolumn{3}{c}{n/a (ORT-Web 1.17 limitation)} \\
\bottomrule
\end{tabular}
\end{table*}


Three remarks. First, \emph{all three students offer a $7$--$19\times$ speedup over the teacher on the M2~Pro target}, primarily reflecting the parameter and FLOP reductions but also benefiting from ORT's CPU-graph optimisations on smaller models. Second, \emph{INT8 yields a further $1.3$--$1.6\times$ speedup on top of the FP32 student baseline}, consistent with INT8 SIMD-CPU expectations. The combined teacher-FP32 $\rightarrow$ EfficientNet-Lite0-INT8 speedup is $15.8\times$ at a 95\% IoU recovery. Third, \emph{the WASM column is roughly $5\times$ slower than the M2-CPU column for the same FP32 model}, which is the expected single-threaded-WebAssembly-vs-native-CPU gap and represents an effective lower bound on browser-deployment performance. INT8 WASM rows are unsupported by ONNX Runtime Web 1.17 at time of writing and are marked n/a.

\subsection{Pareto Analysis}
\label{sec:results-pareto}

\begin{figure}[ht]
\centering
\includegraphics[width=0.95\linewidth]{figures/F2_pareto.pdf}
\caption{\textbf{Pareto frontier on four deployment-relevant axes.} IoU vs.~$\{$Params, FP32 ONNX size, M2~Pro p50, WASM p50$\}$. The \textbf{EfficientNet-Lite0 INT8} configuration is on the Pareto frontier on the parameter, M2-latency, and WASM-latency axes simultaneously; \textbf{MobileViT-XXS no-KD FP32} is Pareto-optimal on the accuracy axis but not deployable as INT8 under the current static QDQ recipe.}
\label{fig:f2-pareto}
\end{figure}

The Pareto front therefore identifies two distinct deployment operating points:

\begin{itemize}
\item \textbf{Maximum-accuracy operating point:} MobileViT-XXS no-KD FP32 (IoU 0.937, 28\,ms M2, 149\,ms WASM, 12\,MB).
\item \textbf{Maximum-throughput-with-acceptable-accuracy operating point:} EfficientNet-Lite0 INT8 (IoU 0.885, 17\,ms M2, no WASM, 5.2\,MB).
\end{itemize}

We recommend the latter for unattended flood-mapping pipelines ($1.6\times$ throughput, $2.3\times$ smaller binary) and the former for human-in-the-loop applications where the 5.2-point IoU drop matters.

\begin{figure}[ht]
\centering
\includegraphics[width=0.95\linewidth]{figures/F4_throughput.png}
\caption{\textbf{Per-platform throughput, FP32 vs.~INT8.} Frames-per-second on M2~Pro (left) and browser WASM SIMD (right); INT8 entries are unsupported on ORT-Web 1.17.}
\label{fig:f4-throughput}
\end{figure}

\subsection{Pipeline Diagram and Qualitative Examples}
\label{sec:results-pipeline}

\begin{figure}[ht]
\centering
% Placeholder until F1_pipeline.pdf is finalised (PowerPoint / draw.io export).
\fbox{\parbox{0.85\linewidth}{\centering \emph{Figure placeholder.} F1\_pipeline.pdf is generated manually (PowerPoint or draw.io) and dropped into \texttt{manuscript/figures/}. Layout: $\text{train teacher} \rightarrow \text{distill students (ablation)} \rightarrow \text{quantize INT8} \rightarrow \text{benchmark on M2 CPU + WASM}$.\\[4pt]}}
\caption{\textbf{FloodLite pipeline.} Train teacher (UNet+EffB0) $\rightarrow$ train students (3 backbones $\times$ 4 KD configs) $\rightarrow$ INT8-quantize the no-KD students via ONNX static QDQ $\rightarrow$ benchmark FP32 and INT8 on M2~Pro CPU and browser WASM.}
\label{fig:f1-pipeline}
\end{figure}

\begin{figure}[ht]
\centering
% Placeholder until F3_qualitative.pdf is rendered by scripts/make_qualitative.py.
\fbox{\parbox{0.85\linewidth}{\centering \emph{Figure placeholder.} F3\_qualitative.pdf is rendered by \texttt{scripts/make\_qualitative.py} once FSSD is on the rendering host; supplementary material.\\[4pt]}}
\caption{\textbf{Qualitative panel} (supplementary). Two rows $\times$ five columns: input, GT, teacher FP32 prediction, EfficientNet-Lite0 INT8 prediction, error map (TP/FP/FN). The easy row is the highest-IoU fold-0 validation image under the deployable model; the hard row is the lowest-IoU image with $\geq 2\%$ flood-pixel coverage.}
\label{fig:f3-qualitative}
\end{figure}

Two persistent failure modes shared by teacher and student are visible in the hard row of Figure~\ref{fig:f3-qualitative}: wet asphalt confused with water, and under-segmentation of vegetation-occluded flood edges. We return to these in Section~\ref{sec:discussion-failures}.
